\documentclass[11pt,a4paper]{article}
\usepackage[utf8]{inputenc}
\usepackage[T1]{fontenc}
\usepackage[catalan]{babel}
\usepackage{lmodern}
\usepackage[sfdefault,lining]{FiraSans}
\usepackage[a4paper,top=2.3cm,bottom=2.2cm,left=1.8cm,right=1.8cm,headheight=1.3cm,headsep=0.4cm,footskip=1cm]{geometry}
\usepackage{amsmath,amssymb}
\usepackage{icomma}
\usepackage{graphicx}
\usepackage[table]{xcolor}
\usepackage{tikz}
\usepackage{fontawesome5}
\usepackage{ccicons}
\usepackage{fancyhdr}
\usepackage{lastpage}
\usepackage{tasks}
\usepackage{needspace}
\usepackage[most]{tcolorbox}
\usepackage{array}
\usepackage[hidelinks,pdftitle={Fitxa d'exercicis de la unitat 5: Trigonometria. 4t d'ESO Matemàtiques B},pdfsubject={Matemàtiques B 4t ESO},pdfkeywords={trigonometria, radians, sinus, cosinus, tangent, relacions fonamentals, triangles rectangles, circumferència goniomètrica, reducció al primer quadrant}]{hyperref}

\definecolor{navy}{HTML}{1B2748}
\definecolor{accent}{HTML}{E8603A}
\definecolor{teal}{HTML}{1F9E90}
\definecolor{soft}{HTML}{EEF1F8}
\definecolor{muted}{HTML}{5D6685}

\setlength{\parindent}{0pt}
\setlength{\parskip}{3pt}
\settasks{label=\alph*), label-format=\bfseries\color{navy}, label-width=1.3em, item-indent=1.8em, column-sep=0.6em, after-item-skip=5pt, before-skip=3pt, after-skip=4pt}

% ---------- capçalera i peu ----------
\newcommand{\llicencia}{\href{https://creativecommons.org/licenses/by-nc-sa/4.0/deed.ca}{\ccbyncsa}\ \ Llicència Creative Commons Reconeixement-NoComercial-CompartirIgual 4.0 Internacional (CC BY-NC-SA 4.0)}
\pagestyle{fancy}
\fancyhf{}
\fancyhead[L]{\raisebox{-0.35\height}{\includegraphics[height=0.95cm]{logo-gva.png}}}
\fancyhead[C]{\small\color{muted}Fitxa d'exercicis · Unitat 5: Trigonometria · 4t ESO B}
\fancyhead[R]{\raisebox{-0.35\height}{\includegraphics[height=0.95cm]{logo-vermellar.png}}}
\fancyfoot[L]{\scriptsize\color{muted}\llicencia}
\fancyfoot[R]{\small\color{muted}\thepage}
\renewcommand{\headrulewidth}{0.4pt}
\renewcommand{\headrule}{\color{muted!40}\hrule height \headrulewidth}
\fancypagestyle{primera}{\fancyhf{}\renewcommand{\headrulewidth}{0pt}
  \fancyfoot[L]{\scriptsize\color{muted}\llicencia}
  \fancyfoot[R]{\small\color{muted}\thepage}}

% ---------- nivells i exercicis ----------
\newcommand{\estrelles}[1]{{\color{accent}\ifcase#1\or\faStar\faStar[regular]\faStar[regular]\or\faStar\faStar\faStar[regular]\or\faStar\faStar\faStar\fi}}
\newcounter{exer}
\newcommand{\ex}[1]{\par\needspace{4\baselineskip}\medskip\refstepcounter{exer}\noindent{\color{navy}\bfseries\theexer.}\ \ {\small\estrelles{#1}}\ \ }

\newcommand{\seccio}[2]{\par\needspace{16\baselineskip}\bigskip
  \begin{tcolorbox}[colback=navy,colframe=navy,coltext=white,arc=2mm,boxrule=0pt,left=3mm,right=3mm,top=1.2mm,bottom=1.2mm]
  {\large\bfseries #1\quad #2}
  \end{tcolorbox}\vspace{-2pt}}

\newcommand{\exf}[1]{\par\medskip\refstepcounter{exer}\noindent{\color{navy}\bfseries\theexer.}\ \ {\small\estrelles{#1}}\ \ }
\newcolumntype{C}[1]{>{\centering\arraybackslash}p{#1}}
% quadre de recordatori (només a la fitxa d'exercicis)
\newcommand{\recorda}[1]{\par\smallskip\begin{tcolorbox}[colback=accent!7,colframe=accent!55,arc=2mm,boxrule=0.5pt,left=2.5mm,right=2.5mm,top=1mm,bottom=1mm]\small{\color{accent}\faLightbulb[regular]}\ \textbf{Recorda:}\ #1\end{tcolorbox}}
\usetikzlibrary{calc}
\usepackage{enumitem}
\providecommand{\tg}{\operatorname{tg}}
\newcommand{\res}[1]{#1}\newcommand{\rest}[1]{#1}

\def\fmneed{11}% que el títol d'una secció no quede sol al final del full
% =====================================================================
% Capçalera FentMates de les fitxes: un color per a cada curs.
% Els colors són els de la pàgina inicial (LEVELS a assets/js/materials.js);
% el logo ve de _eines/logo/pdf/ (python3 _eines/logo/genera.py).
%
% Ús, al final del preàmbul de la unitat:
%   % =====================================================================
% Capçalera FentMates de les fitxes: un color per a cada curs.
% Els colors són els de la pàgina inicial (LEVELS a assets/js/materials.js);
% el logo ve de _eines/logo/pdf/ (python3 _eines/logo/genera.py).
%
% Ús, al final del preàmbul de la unitat:
%   % =====================================================================
% Capçalera FentMates de les fitxes: un color per a cada curs.
% Els colors són els de la pàgina inicial (LEVELS a assets/js/materials.js);
% el logo ve de _eines/logo/pdf/ (python3 _eines/logo/genera.py).
%
% Ús, al final del preàmbul de la unitat:
%   \input{../comu/fentmates.tex}
%   \fmcurs{1eso}
% i, al començament del document:
%   \capcalera{Matemàtiques}{1r ESO}{Fitxa d'exercicis · Unitat 2}{Divisibilitat}
% Opcional abans de \input: \def\fmneed{14} (espai mínim abans de cada secció).
% =====================================================================
\makeatletter
\def\fm@logodir{../../logo/pdf}
\providecommand{\fmneed}{10}
\def\fm@cursnom{}

% clau del curs, color i color del text damunt del color
\newcommand{\fm@curs}[3]{\@namedef{fm@color@#1}{#2}\@namedef{fm@text@#1}{#3}}
\fm@curs{1eso}{E6B84A}{navy}
\fm@curs{2eso}{3B82F6}{white}
\fm@curs{3eso}{FF7A50}{white}
\fm@curs{4eso-a}{A78BFA}{white}
\fm@curs{4eso-b}{6D28D9}{white}
\fm@curs{1batx-mi}{9B1C1C}{white}
\fm@curs{1batx-ccss}{EF6B6B}{white}
\fm@curs{2batx-mii}{1E6B47}{white}
\fm@curs{2batx-ccss}{2BB3A3}{white}

\newcommand{\fmcurs}[1]{%
  \@ifundefined{fm@color@#1}{\PackageError{fentmates}{Curs desconegut: #1}{}}{}%
  \edef\fm@tmp{\noexpand\definecolor{curs}{HTML}{\@nameuse{fm@color@#1}}%
               \noexpand\colorlet{curstext}{\@nameuse{fm@text@#1}}}\fm@tmp
  \colorlet{cursfosc}{curs!72!black}%
  \def\fm@logo{\fm@logodir/vertical-#1.pdf}%
  \def\fm@marca{\fm@logodir/marca-#1.pdf}}

% ---------- capçalera de la primera pàgina ----------
% #1 matèria (petit), #2 curs, #3 tipus de document i unitat, #4 títol
\newcommand{\capcalera}[4]{%
  \gdef\fm@cursnom{#2}%
  \par\noindent
  \begin{tcolorbox}[colback=curs,colframe=curs,boxrule=0pt,arc=2.2mm,width=3.5cm,height=2.15cm,
      valign=center,left=3mm,right=2mm,top=1mm,bottom=1mm,nobeforeafter,box align=center]
    \color{curstext}{\footnotesize #1}\par\vspace{1pt}{\fontsize{19}{21}\selectfont\bfseries #2\par}%
  \end{tcolorbox}\hspace{5mm}%
  \begin{minipage}[c]{\dimexpr\linewidth-3.5cm-5mm-2.9cm\relax}
    {\color{cursfosc}\small\bfseries\MakeUppercase{#3}}\par\vspace{2pt}
    {\color{navy}\fontsize{20}{23}\selectfont\bfseries #4\par}
  \end{minipage}\hfill
  \begin{minipage}[c]{2.6cm}\raggedleft\includegraphics[width=2.6cm]{\fm@logo}\end{minipage}\par
  \vspace{2pt}}

% ---------- barres de secció del color del curs ----------
\renewcommand{\seccio}[2]{\par\needspace{\fmneed\baselineskip}\bigskip
  \begin{tcolorbox}[colback=curs,colframe=curs,coltext=curstext,arc=2mm,boxrule=0pt,left=3mm,right=3mm,top=1.2mm,bottom=1.2mm]
  {\large\bfseries #1\quad #2}
  \end{tcolorbox}\vspace{-2pt}}

% ---------- estrelles de dificultat del color del curs (a Batxillerat no se'n posen) ----------
\@ifundefined{estrelles}{}{%
  \renewcommand{\estrelles}[1]{{\color{curs}\ifcase#1\or\faStar\faStar[regular]\faStar[regular]\or\faStar\faStar\faStar[regular]\or\faStar\faStar\faStar\fi}}}

% ---------- capçalera i peu de les pàgines ----------
% primera pàgina: logos institucionals al peu
\fancypagestyle{primera}{\fancyhf{}\renewcommand{\headrulewidth}{0pt}
  \fancyfoot[L]{\raisebox{-0.32\height}{\includegraphics[height=0.78cm]{logo-gva.png}}\hspace{5mm}%
    \raisebox{-0.32\height}{\includegraphics[height=0.78cm]{logo-vermellar.png}}}
  \fancyfoot[R]{\scriptsize\color{muted}\href{https://creativecommons.org/licenses/by-nc-sa/4.0/deed.ca}{\ccbyncsa}\ \ CC BY-NC-SA 4.0 · fentmates.org\qquad\small\thepage}}
% la resta: el símbol a l'esquerra i el curs a la dreta, amb el títol al centre com abans
\fancyhead[L]{\raisebox{-0.3\height}{\includegraphics[height=0.62cm]{\fm@marca}}}
\fancyhead[R]{\small\bfseries\color{cursfosc}\fm@cursnom}
\renewcommand{\headrule}{\color{curs!45}\hrule height \headrulewidth}
\makeatother

%   \fmcurs{1eso}
% i, al començament del document:
%   \capcalera{Matemàtiques}{1r ESO}{Fitxa d'exercicis · Unitat 2}{Divisibilitat}
% Opcional abans de \input: \def\fmneed{14} (espai mínim abans de cada secció).
% =====================================================================
\makeatletter
\def\fm@logodir{../../logo/pdf}
\providecommand{\fmneed}{10}
\def\fm@cursnom{}

% clau del curs, color i color del text damunt del color
\newcommand{\fm@curs}[3]{\@namedef{fm@color@#1}{#2}\@namedef{fm@text@#1}{#3}}
\fm@curs{1eso}{E6B84A}{navy}
\fm@curs{2eso}{3B82F6}{white}
\fm@curs{3eso}{FF7A50}{white}
\fm@curs{4eso-a}{A78BFA}{white}
\fm@curs{4eso-b}{6D28D9}{white}
\fm@curs{1batx-mi}{9B1C1C}{white}
\fm@curs{1batx-ccss}{EF6B6B}{white}
\fm@curs{2batx-mii}{1E6B47}{white}
\fm@curs{2batx-ccss}{2BB3A3}{white}

\newcommand{\fmcurs}[1]{%
  \@ifundefined{fm@color@#1}{\PackageError{fentmates}{Curs desconegut: #1}{}}{}%
  \edef\fm@tmp{\noexpand\definecolor{curs}{HTML}{\@nameuse{fm@color@#1}}%
               \noexpand\colorlet{curstext}{\@nameuse{fm@text@#1}}}\fm@tmp
  \colorlet{cursfosc}{curs!72!black}%
  \def\fm@logo{\fm@logodir/vertical-#1.pdf}%
  \def\fm@marca{\fm@logodir/marca-#1.pdf}}

% ---------- capçalera de la primera pàgina ----------
% #1 matèria (petit), #2 curs, #3 tipus de document i unitat, #4 títol
\newcommand{\capcalera}[4]{%
  \gdef\fm@cursnom{#2}%
  \par\noindent
  \begin{tcolorbox}[colback=curs,colframe=curs,boxrule=0pt,arc=2.2mm,width=3.5cm,height=2.15cm,
      valign=center,left=3mm,right=2mm,top=1mm,bottom=1mm,nobeforeafter,box align=center]
    \color{curstext}{\footnotesize #1}\par\vspace{1pt}{\fontsize{19}{21}\selectfont\bfseries #2\par}%
  \end{tcolorbox}\hspace{5mm}%
  \begin{minipage}[c]{\dimexpr\linewidth-3.5cm-5mm-2.9cm\relax}
    {\color{cursfosc}\small\bfseries\MakeUppercase{#3}}\par\vspace{2pt}
    {\color{navy}\fontsize{20}{23}\selectfont\bfseries #4\par}
  \end{minipage}\hfill
  \begin{minipage}[c]{2.6cm}\raggedleft\includegraphics[width=2.6cm]{\fm@logo}\end{minipage}\par
  \vspace{2pt}}

% ---------- barres de secció del color del curs ----------
\renewcommand{\seccio}[2]{\par\needspace{\fmneed\baselineskip}\bigskip
  \begin{tcolorbox}[colback=curs,colframe=curs,coltext=curstext,arc=2mm,boxrule=0pt,left=3mm,right=3mm,top=1.2mm,bottom=1.2mm]
  {\large\bfseries #1\quad #2}
  \end{tcolorbox}\vspace{-2pt}}

% ---------- estrelles de dificultat del color del curs (a Batxillerat no se'n posen) ----------
\@ifundefined{estrelles}{}{%
  \renewcommand{\estrelles}[1]{{\color{curs}\ifcase#1\or\faStar\faStar[regular]\faStar[regular]\or\faStar\faStar\faStar[regular]\or\faStar\faStar\faStar\fi}}}

% ---------- capçalera i peu de les pàgines ----------
% primera pàgina: logos institucionals al peu
\fancypagestyle{primera}{\fancyhf{}\renewcommand{\headrulewidth}{0pt}
  \fancyfoot[L]{\raisebox{-0.32\height}{\includegraphics[height=0.78cm]{logo-gva.png}}\hspace{5mm}%
    \raisebox{-0.32\height}{\includegraphics[height=0.78cm]{logo-vermellar.png}}}
  \fancyfoot[R]{\scriptsize\color{muted}\href{https://creativecommons.org/licenses/by-nc-sa/4.0/deed.ca}{\ccbyncsa}\ \ CC BY-NC-SA 4.0 · fentmates.org\qquad\small\thepage}}
% la resta: el símbol a l'esquerra i el curs a la dreta, amb el títol al centre com abans
\fancyhead[L]{\raisebox{-0.3\height}{\includegraphics[height=0.62cm]{\fm@marca}}}
\fancyhead[R]{\small\bfseries\color{cursfosc}\fm@cursnom}
\renewcommand{\headrule}{\color{curs!45}\hrule height \headrulewidth}
\makeatother

%   \fmcurs{1eso}
% i, al començament del document:
%   \capcalera{Matemàtiques}{1r ESO}{Fitxa d'exercicis · Unitat 2}{Divisibilitat}
% Opcional abans de \input: \def\fmneed{14} (espai mínim abans de cada secció).
% =====================================================================
\makeatletter
\def\fm@logodir{../../logo/pdf}
\providecommand{\fmneed}{10}
\def\fm@cursnom{}

% clau del curs, color i color del text damunt del color
\newcommand{\fm@curs}[3]{\@namedef{fm@color@#1}{#2}\@namedef{fm@text@#1}{#3}}
\fm@curs{1eso}{E6B84A}{navy}
\fm@curs{2eso}{3B82F6}{white}
\fm@curs{3eso}{FF7A50}{white}
\fm@curs{4eso-a}{A78BFA}{white}
\fm@curs{4eso-b}{6D28D9}{white}
\fm@curs{1batx-mi}{9B1C1C}{white}
\fm@curs{1batx-ccss}{EF6B6B}{white}
\fm@curs{2batx-mii}{1E6B47}{white}
\fm@curs{2batx-ccss}{2BB3A3}{white}

\newcommand{\fmcurs}[1]{%
  \@ifundefined{fm@color@#1}{\PackageError{fentmates}{Curs desconegut: #1}{}}{}%
  \edef\fm@tmp{\noexpand\definecolor{curs}{HTML}{\@nameuse{fm@color@#1}}%
               \noexpand\colorlet{curstext}{\@nameuse{fm@text@#1}}}\fm@tmp
  \colorlet{cursfosc}{curs!72!black}%
  \def\fm@logo{\fm@logodir/vertical-#1.pdf}%
  \def\fm@marca{\fm@logodir/marca-#1.pdf}}

% ---------- capçalera de la primera pàgina ----------
% #1 matèria (petit), #2 curs, #3 tipus de document i unitat, #4 títol
\newcommand{\capcalera}[4]{%
  \gdef\fm@cursnom{#2}%
  \par\noindent
  \begin{tcolorbox}[colback=curs,colframe=curs,boxrule=0pt,arc=2.2mm,width=3.5cm,height=2.15cm,
      valign=center,left=3mm,right=2mm,top=1mm,bottom=1mm,nobeforeafter,box align=center]
    \color{curstext}{\footnotesize #1}\par\vspace{1pt}{\fontsize{19}{21}\selectfont\bfseries #2\par}%
  \end{tcolorbox}\hspace{5mm}%
  \begin{minipage}[c]{\dimexpr\linewidth-3.5cm-5mm-2.9cm\relax}
    {\color{cursfosc}\small\bfseries\MakeUppercase{#3}}\par\vspace{2pt}
    {\color{navy}\fontsize{20}{23}\selectfont\bfseries #4\par}
  \end{minipage}\hfill
  \begin{minipage}[c]{2.6cm}\raggedleft\includegraphics[width=2.6cm]{\fm@logo}\end{minipage}\par
  \vspace{2pt}}

% ---------- barres de secció del color del curs ----------
\renewcommand{\seccio}[2]{\par\needspace{\fmneed\baselineskip}\bigskip
  \begin{tcolorbox}[colback=curs,colframe=curs,coltext=curstext,arc=2mm,boxrule=0pt,left=3mm,right=3mm,top=1.2mm,bottom=1.2mm]
  {\large\bfseries #1\quad #2}
  \end{tcolorbox}\vspace{-2pt}}

% ---------- estrelles de dificultat del color del curs (a Batxillerat no se'n posen) ----------
\@ifundefined{estrelles}{}{%
  \renewcommand{\estrelles}[1]{{\color{curs}\ifcase#1\or\faStar\faStar[regular]\faStar[regular]\or\faStar\faStar\faStar[regular]\or\faStar\faStar\faStar\fi}}}

% ---------- capçalera i peu de les pàgines ----------
% primera pàgina: logos institucionals al peu
\fancypagestyle{primera}{\fancyhf{}\renewcommand{\headrulewidth}{0pt}
  \fancyfoot[L]{\raisebox{-0.32\height}{\includegraphics[height=0.78cm]{logo-gva.png}}\hspace{5mm}%
    \raisebox{-0.32\height}{\includegraphics[height=0.78cm]{logo-vermellar.png}}}
  \fancyfoot[R]{\scriptsize\color{muted}\href{https://creativecommons.org/licenses/by-nc-sa/4.0/deed.ca}{\ccbyncsa}\ \ CC BY-NC-SA 4.0 · fentmates.org\qquad\small\thepage}}
% la resta: el símbol a l'esquerra i el curs a la dreta, amb el títol al centre com abans
\fancyhead[L]{\raisebox{-0.3\height}{\includegraphics[height=0.62cm]{\fm@marca}}}
\fancyhead[R]{\small\bfseries\color{cursfosc}\fm@cursnom}
\renewcommand{\headrule}{\color{curs!45}\hrule height \headrulewidth}
\makeatother

\fmcurs{4eso-b}

\begin{document}
\thispagestyle{primera}

\capcalera{Matemàtiques B}{4t ESO}{Fitxa d'exercicis · Unitat 5}{\raggedright Trigonometria}
\vspace{10pt}
Nom i cognoms: \rule{8.4cm}{0.4pt}\hfill Grup: \rule{1.6cm}{0.4pt}\hfill Data: \rule{2.2cm}{0.4pt}
\vspace{8pt}

\begin{tcolorbox}[colback=soft,colframe=soft,arc=2mm,left=3mm,right=3mm,top=1.5mm,bottom=1.5mm]
\small Exercicis graduats per nivells de dificultat:\quad \estrelles{1}\ bàsic \qquad \estrelles{2}\ mitjà \qquad \estrelles{3}\ avançat.\\[2pt] Resol-los al quadern mostrant tots els passos i fent un dibuix en els problemes. Fes servir la calculadora en mode \texttt{D} (graus) i dona els angles en graus, minuts i segons i les longituds amb dos decimals, tret que l'enunciat demane el resultat exacte. La fitxa recull els exercicis de les fitxes del departament i n'afegeix de nous.
\end{tcolorbox}

% =====================================================
\seccio{1}{Mesura d'angles: graus i radians}

\exf{1} Expressa en radians aquests angles (de manera exacta, en funció de $\pi$):\begin{tasks}(6)
  \task $45^\circ$
  \task $30^\circ$
  \task $120^\circ$
  \task $150^\circ$
  \task $540^\circ$
  \task $1080^\circ$
  \task $1800^\circ$
  \task $240^\circ$
  \task $310^\circ$
  \task $70^\circ$
  \task $12^\circ$
  \task $720^\circ$
\end{tasks}

\ex{1} Expressa en graus sexagesimals aquests angles en radians:\begin{tasks}(6)
  \task $\pi$
  \task $4\pi$
  \task $7\pi$
  \task $\frac{\pi}{6}$
  \task $\frac{3\pi}{4}$
  \task $\frac{3\pi}{2}$
  \task $\frac{5\pi}{3}$
  \task $\frac{\pi}{2}$
  \task $\frac{4\pi}{3}$
  \task $1{,}2\pi$
  \task $\frac{2\pi}{5}$
  \task $\frac{\pi}{9}$
\end{tasks}

\ex{1} Completa la taula següent:\begin{center}\renewcommand{\arraystretch}{1.6}\begin{tabular}{|l|*{10}{C{1.05cm}|}}\hline\textbf{Graus} & $30^\circ$ &  & $60^\circ$ &  & $135^\circ$ &  & $210^\circ$ &  & $300^\circ$ & \\\hline\textbf{Radians} &  & $\frac{\pi}{4}$ &  & $\frac{2\pi}{3}$ &  & $\frac{5\pi}{6}$ &  & $\frac{5\pi}{4}$ &  & $\frac{11\pi}{6}$\\\hline\end{tabular}\end{center}

\ex{1} Expressa, tant en graus sexagesimals com en radians, aquestes parts d'una circumferència:\begin{tasks}(4)
  \task Una quarta part.
  \task Tres quartes parts.
  \task Una sisena part.
  \task Una tercera part.
\end{tasks}

\ex{2} En l'hexàgon regular $ABCDEF$ de centre $O$, expressa en radians els angles $\widehat{AOB}$, $\widehat{ABC}$, $\widehat{BAC}$ i $\widehat{ACD}$.\begin{center}\begin{tikzpicture}[scale=1.25]
\foreach \i/\n in {0/D,1/E,2/F,3/A,4/B,5/C} \coordinate (\n) at (60*\i:1.2);
\coordinate (O) at (0,0);
\draw[navy, line width=1.2pt] (A)--(B)--(C)--(D)--(E)--(F)--cycle;
\draw[muted, dashed] (A)--(C) (A)--(D) (O)--(B);
\fill (O) circle (1.3pt) node[above right=-1pt,font=\small]{$O$};
\foreach \n/\p in {A/left,B/below left,C/below right,D/right,E/above right,F/above left} \node[\p,font=\small] at (\n) {$\n$};
\end{tikzpicture}\end{center}

\ex{2} Passa a graus, minuts i segons: \ $1$ rad, \ $0{,}75$ rad \ i \ $2{,}5$ rad. Passa a radians, de manera exacta: \ $22^\circ30'$, \ $52^\circ30'$ \ i \ $157^\circ30'$.

\ex{3} La busca de les hores d'un rellotge fa una volta en 12 hores. Quin angle, en graus i en radians, recorre en una hora? I la busca dels minuts en 20 minuts? Quin angle formen les dues busques a les 4 en punt?


% =====================================================
\seccio{2}{Raons trigonomètriques d'un angle agut}

\exf{1} Calcula les raons trigonomètriques (sinus, cosinus i tangent) de l'angle agut $\alpha$ i de l'altre angle agut, $\beta$, d'aquests triangles rectangles:\begin{enumerate}[label=\alph*),leftmargin=*,itemsep=2pt]\item Catets de 30 cm i 16 cm; $\alpha$ és oposat al catet de 16 cm.\item Hipotenusa de 10 cm i catet de 6 cm oposat a $\alpha$.\item Catets de 24 cm i 10 cm; $\alpha$ és oposat al catet de 24 cm.\item Hipotenusa de 25 cm i catet de 7 cm contigu a $\alpha$.\item Hipotenusa de 29 cm i catet de 20 cm contigu a $\alpha$.\end{enumerate}

\ex{1} Calcula les raons trigonomètriques dels angles $\hat A$, $\hat C$, $\widehat{ABD}$ i $\widehat{CBD}$.\begin{center}\begin{tikzpicture}[scale=0.27]
\coordinate (A) at (0,0); \coordinate (D) at (9,0); \coordinate (B) at (9,12); \coordinate (C) at (25,0);
\draw[navy, line width=1.1pt] (A)--(B)--(C)--cycle; \draw[muted, dashed] (B)--(D);
\draw (D)++(-0.9,0)--++(0,0.9)--++(0.9,0);
\node[below left] at (A) {$A$}; \node[above] at (B) {$B$}; \node[below right] at (C) {$C$}; \node[below] at (D) {$D$};
\node[above left, font=\small] at (4.5,6) {15 cm}; \node[right, font=\small] at (9,5) {12 cm}; \node[below, font=\small] at (17,0) {16 cm};
\end{tikzpicture}\end{center}

\ex{1} Completa la taula amb la calculadora (amb quatre decimals):\begin{center}\renewcommand{\arraystretch}{1.45}\begin{tabular}{|c|*{4}{C{2.3cm}|}}\hline $\alpha$ & $15^\circ$ & $55^\circ20'$ & $72^\circ25'40''$ & $85{,}5^\circ$\\\hline
$\sin\alpha$ & & & & \\\hline $\cos\alpha$ & & & & \\\hline $\tg\alpha$ & & & & \\\hline\end{tabular}\end{center}

\ex{1} Troba, amb la calculadora, l'angle agut $\alpha$ en cada cas:\begin{tasks}(5)
  \task $\sin\alpha=\frac{\sqrt3}{2}$
  \task $\cos\alpha=\frac13$
  \task $\tg\alpha=2{,}5$
  \task $\sin\alpha=0{,}25$
  \task $\cos\alpha=0{,}8$
\end{tasks}

\ex{2} Comprova, per a $\alpha=60^\circ$, a partir dels hexàgons regulars de costat 4 cm i de costat 6 cm (formats per sis triangles equilàters), que les raons trigonomètriques no varien segons la longitud del costat.

\ex{2} Calcula els costats que falten en cada triangle rectangle i, si cal, les altres raons de l'angle:\begin{enumerate}[label=\alph*),leftmargin=*,itemsep=2pt]\item El catet contigu a un angle $\alpha$ mesura 5 cm i $\tg\alpha=\frac65$.\item El sinus d'un angle val $\frac{12}{13}$ i el catet oposat mesura 24 cm. Calcula també el cosinus i la tangent.\item El catet oposat a un angle mesura 15 cm i la tangent d'aquest angle val $\frac34$. Calcula també el sinus i el cosinus.\item La hipotenusa mesura 2 cm i la tangent d'un dels angles val 1.\end{enumerate}

\ex{3} La hipotenusa d'un triangle rectangle mesura 10 cm i la tangent d'un dels angles aguts val $\frac{7}{12}$. Calcula els catets i el sinus i el cosinus d'aquest angle.


% =====================================================
\seccio{3}{Raons trigonomètriques de 30°, 45° i 60°}

\exf{1} Completa la taula següent sense calculadora, de manera exacta:\begin{center}\renewcommand{\arraystretch}{1.6}\begin{tabular}{|l|*{5}{C{1.6cm}|}}\hline
& $0^\circ$ & $30^\circ$ & $45^\circ$ & $60^\circ$ & $90^\circ$\\\hline
$\sin$ & & & & & \\\hline $\cos$ & & & & & \\\hline $\tg$ & & & & & \\\hline\end{tabular}\end{center}

\ex{1} En cada triangle rectangle es coneixen un angle agut i un costat. Calcula els altres dos costats de manera exacta:\begin{enumerate}[label=\alph*),leftmargin=*,itemsep=2pt]\item Un angle de $60^\circ$ i el catet oposat, de 5 cm.\item Un angle de $30^\circ$ i el catet oposat, de 2 cm.\item Un angle de $45^\circ$ i el catet contigu, de 7 cm.\item Un angle de $60^\circ$ i el catet contigu, de 6 cm.\item Un angle de $45^\circ$ i el catet oposat, de 12 cm.\item Un angle de $30^\circ$ i la hipotenusa, de 30 cm.\end{enumerate}

\ex{1} Calcula sense calculadora:\begin{tasks}(3)
  \task $\sin30^\circ+\cos60^\circ$
  \task $\tg45^\circ-2\sin^245^\circ$
  \task $\sin60^\circ\cos30^\circ-\sin30^\circ\cos60^\circ$
  \task $\tg60^\circ\cdot\tg30^\circ$
  \task $(\sin45^\circ+\cos45^\circ)^2$
  \task $\cos^230^\circ-\sin^230^\circ$
\end{tasks}

\ex{2} Sense fer servir el teorema de Pitàgores, calcula l'altura d'un triangle equilàter de 8 cm de costat i la diagonal d'un quadrat de 5 cm de costat.

\ex{2} Un hexàgon regular es pot dividir en 6 triangles equilàters. Si el costat de l'hexàgon mesura 4 cm, calcula'n l'apotema i l'àrea de manera exacta.

\ex{3} Dues antenes de ràdio, de 100 m i 75 m, estan subjectes al terra com s'indica en la figura. Calcula la longitud de cada tram de cable i la distància $\overline{AE}$. Dona els resultats de manera exacta.\begin{center}\begin{tikzpicture}[scale=0.62]
\coordinate (A) at (0,0); \coordinate (P) at (2.3,0); \coordinate (B) at (2.3,4); \coordinate (C) at (9.2,0); \coordinate (Q) at (12.2,0); \coordinate (D) at (12.2,3); \coordinate (E) at (13.9,0);
\draw[muted] (-0.6,0)--(14.6,0);
\draw[accent, line width=1.6pt] (P)--(B) (Q)--(D);
\draw[navy, dashed] (A)--(B)--(C)--(D)--(E);
\foreach \n/\p in {A/below,P/below,C/below,Q/below,E/below,B/above,D/above} \node[\p,font=\small] at (\n) {$\n$};
\node[left,font=\small] at (2.3,2) {100 m}; \node[left,font=\small] at (12.2,1.4) {75 m};
\draw (0.8,0) arc (0:60:0.8); \node[font=\scriptsize] at (1.15,0.4) {$60^\circ$};
\draw (8.2,0) arc (180:150:1); \node[font=\scriptsize] at (7.7,0.3) {$30^\circ$};
\draw (12.2,2.2) arc (270:225:0.8); \node[font=\scriptsize] at (11.7,1.9) {$45^\circ$};
\draw (12.2,2.2) arc (270:300:0.8); \node[font=\scriptsize] at (12.75,1.85) {$30^\circ$};
\end{tikzpicture}\end{center}


% =====================================================
\seccio{4}{Relacions fonamentals}

\exf{1} Sabent que l'angle $\alpha$ és agut, calcula'n, en cada cas, les altres raons trigonomètriques de manera exacta:\begin{tasks}(5)
  \task $\sin\alpha=\frac{7}{25}$
  \task $\sin\alpha=\frac{12}{37}$
  \task $\sin\alpha=\frac{11}{61}$
  \task $\cos\alpha=\frac{20}{29}$
  \task $\cos\alpha=\frac{12}{13}$
  \task $\sin\alpha=\frac{3}{8}$
  \task $\cos\alpha=\frac{1}{3}$
  \task $\sin\alpha=\frac{15}{34}$
  \task $\cos\alpha=\frac{2}{5}$
  \task $\sin\alpha=\frac{4}{7}$
\end{tasks}\recorda{$\sin^2\alpha+\cos^2\alpha=1$, \quad $\tg\alpha=\dfrac{\sin\alpha}{\cos\alpha}$, \quad $\tg^2\alpha+1=\dfrac{1}{\cos^2\alpha}$.}

\ex{2} Si l'angle $\alpha$ és agut, calcula, en cada cas, el sinus i el cosinus de l'angle de manera exacta:\begin{tasks}(7)
  \task $\tg\alpha=\frac{5}{12}$
  \task $\tg\alpha=\frac{3}{4}$
  \task $\tg\alpha=\frac{5}{3}$
  \task $\tg\alpha=\frac{5}{13}$
  \task $\tg\alpha=\frac{5}{8}$
  \task $\tg\alpha=3$
  \task $\tg\alpha=1$
\end{tasks}

\ex{1} Sabent que l'angle $\alpha$ és agut, calcula les altres dues raons (amb quatre decimals):\begin{tasks}(3)
  \task $\cos\alpha=0,1357$
  \task $\cos\alpha=0,9531$
  \task $\cos\alpha=0,2864$
  \task $\sin\alpha=0,1827$
  \task $\sin\alpha=0,9542$
  \task $\sin\alpha=0,4531$
\end{tasks}

\ex{2} Comprova si hi ha algun angle que tinga aquestes raons trigonomètriques:\begin{tasks}(2)
  \task $\sin\alpha=\frac{7}{25}$ \ i \ $\cos\alpha=\frac{11}{25}$
  \task $\sin\alpha=\frac{7}{13}$ \ i \ $\cos\alpha=\frac{12}{13}$
  \task $\sin\alpha=\frac{\sqrt{5}}{5}$ \ i \ $\cos\alpha=\frac{2 \sqrt{5}}{5}$
  \task $\sin\alpha=\frac{3}{5}$ \ i \ $\cos\alpha=\frac{4}{5}$
  \task $\sin\alpha=\frac{7}{11}$ \ i \ $\tg\alpha=\frac{6}{11}$
  \task $\cos\alpha=\frac{5}{13}$ \ i \ $\tg\alpha=1$
  \task $\cos\alpha=\frac{\sqrt{3}}{2}$ \ i \ $\tg\alpha=\frac{\sqrt{3}}{3}$
  \task $\cos\alpha=\frac{3}{5}$ \ i \ $\tg\alpha=\frac{7}{5}$
\end{tasks}

\ex{2} Comprova si aquests parells de raons trigonomètriques (arrodonides a quatre decimals) poden ser del mateix angle:\begin{tasks}(2)
  \task $\cos\alpha=0,1736$ \ i \ $\tg\alpha=0,4663$
  \task $\cos\alpha=0,9397$ \ i \ $\tg\alpha=0,3640$
  \task $\sin\alpha=0,1485$ \ i \ $\cos\alpha=0,2588$
  \task $\cos\alpha=0,7313$ \ i \ $\tg\alpha=0,9325$
  \task $\sin\alpha=0,2754$ \ i \ $\cos\alpha=0,6691$
\end{tasks}\recorda{amb decimals arrodonits, els dos membres de la relació poden diferir una mica en l'última xifra.}

\ex{3} Fes servir les relacions fonamentals per a demostrar aquestes igualtats:\begin{tasks}(1)
  \task $\dfrac{1}{\tg^2\alpha}+1=\dfrac{1}{\sin^2\alpha}$
  \task $(1+\tg^2\alpha)\cdot\cos^2\alpha=1$
  \task $\sin\alpha\cdot\cos\alpha\cdot\tg\alpha+\cos^2\alpha=1$
\end{tasks}


% =====================================================
\seccio{5}{Resolució de triangles rectangles}

\exf{1} Resol els triangles rectangles següents. En tots, $\hat A=90^\circ$; per tant, $a$ és la hipotenusa.\begin{tasks}(2)
  \task $b=4$ m, \ $\hat B=30^\circ$
  \task $a=7$ m, \ $b=5$ m
  \task $a=10$ m, \ $\hat B=40^\circ$
  \task $c=5$ cm, \ $\sin\hat C=0{,}25$
\end{tasks}

\ex{1} Calcula l'altura $h$, traçada des del vèrtex $B$, de cada triangle $ABC$:\begin{tasks}(2)
  \task $\overline{AB}=18$ cm i $\hat A=62^\circ$.
  \task $\overline{BC}=28$ cm i $\hat C=35^\circ$.
\end{tasks}

\ex{1} Una escala de 3 m està recolzada en una paret. Quin angle forma l'escala amb el terra si la base està a 1,2 m de la paret?

\ex{1} Una torre de 20 m d'altura projecta una ombra de 25 m de llargària. Calcula la inclinació dels raigs solars en aquest moment.

\ex{2} En un triangle isòsceles, el costat desigual mesura 18 m i l'altura, 10 m. Quant mesuren els seus angles?

\ex{2} En una ruta de muntanya, un senyal indica una altitud de 785 m. Tres quilòmetres més endavant (en horitzontal), l'altitud és de 1\,265 m. Troba el pendent mitjà de la ruta i l'angle que forma amb l'horitzontal.

\ex{2} Un avió que és a punt d'aterrar descendeix 2\,160 m en recórrer una distància de 18,41 km. Quin és l'angle $\alpha$ amb què descendeix?

\ex{2} Els braços d'un compàs, que mesuren 12 cm, formen un angle de $50^\circ$. Quin és el radi de la circumferència que es pot traçar amb aquesta obertura?

\ex{2} Dos amics surten del punt $A$. Un es dirigeix al punt $B$ i l'altre al punt $C$, seguint trajectòries rectilínies que formen entre si un angle de $37^\circ$. Si els punts $B$ i $C$ estan separats 120 m i la direcció $BC$ és perpendicular a la direcció $AC$, quants metres recorre un dels amics més que l'altre?


% =====================================================
\seccio{6}{Problemes de trigonometria}

\exf{2} Des d'un cert lloc es veu el punt més alt d'una torre amb un angle d'elevació de $35^\circ$. Si retrocedim 200 m, es veu la mateixa torre amb un angle de $20^\circ$. Calcula l'altura de la torre.

\ex{2} Per calcular l'altura d'un edifici es mesuren els angles d'elevació des de dos punts alineats amb el peu de l'edifici i situats a una distància de 100 m l'un de l'altre. Quina és l'altura de l'edifici si els angles són $33^\circ$ i $46^\circ$?

\ex{2} Per mesurar l'altura d'una muntanya es mesuren els angles d'elevació des de dos punts alineats amb el cim, situats a una distància de 480 m l'un de l'altre i a una altitud de 1\,200 m sobre el nivell de la mar. Quina és l'altura de la muntanya si els angles són $45^\circ$ i $76^\circ$?

\ex{2} Per trobar l'altura a la qual es troba un globus, la Rosa es col·loca en un punt $B$, i jo en un punt $A$, a 5 metres d'ella, de manera que $A$, $B$ i el peu $C$ de la vertical del globus queden alineats ($B$ entre $A$ i $C$). Si des d'$A$ veig el globus amb un angle de $40^\circ$ i la Rosa el veu amb un angle de $50^\circ$, a quina altura es troba el globus?

\ex{2} Un avió vola entre dues ciutats, $A$ i $B$, que estan a 75 km l'una de l'altra. Les visuals des d'$A$ i des de $B$ cap a l'avió formen amb l'horitzontal angles de $36^\circ$ i $12^\circ$, respectivament. Calcula l'altura a la qual vola l'avió i les distàncies a les quals es troba d'$A$ i de $B$ (l'avió i les ciutats estan en el mateix pla vertical).

\ex{2} Volem penjar un llum del sostre d'una habitació amb una corda lligada en dos punts del sostre separats 140 cm, de manera que els angles entre la corda i el sostre siguen de $60^\circ$ i $40^\circ$. Quina serà la longitud total de la corda? A quina distància del sostre quedarà el llum?\begin{center}\begin{tikzpicture}[scale=0.032]
\draw[navy, line width=1.2pt] (-10,0)--(150,0); \foreach \x in {-6,4,...,146} \draw[navy] (\x,0)--++(4,4);
\coordinate (P) at (0,0); \coordinate (Q) at (140,0); \coordinate (L) at (45.69,-79.14);
\draw[muted, line width=0.9pt] (P)--(L)--(Q); \draw[muted, dashed] (L)--(45.69,0);
\fill[accent] (L) circle (2.6);
\draw (12,0) arc (0:-60:12); \node[font=\scriptsize] at (19,-9) {$60^\circ$};
\draw (128,0) arc (180:220:12); \node[font=\scriptsize] at (117,-7) {$40^\circ$};
\draw[<->, muted] (0,14)--(140,14); \node[above, font=\small] at (70,14) {140 cm};
\end{tikzpicture}\end{center}

\ex{2} Una antena de ràdio està subjecta al terra per dos cables d'acer, com indica la figura. Calcula, de manera exacta, l'altura de l'antena, la longitud dels cables i l'angle $\widehat{ABC}$.\begin{center}\begin{tikzpicture}[scale=0.055]
\coordinate (A) at (0,0); \coordinate (C) at (126,0); \coordinate (H) at (31.5,0); \coordinate (B) at (31.5,54.56);
\draw[muted] (-6,0)--(132,0); \draw[navy] (A)--(B)--(C); \draw[accent, line width=1.6pt] (H)--(B);
\node[left] at (A) {$A$}; \node[right] at (C) {$C$}; \node[above] at (B) {$B$};
\draw (9,0) arc (0:60:9); \node[font=\scriptsize] at (13,6) {$60^\circ$};
\draw (114,0) arc (180:150:12); \node[font=\scriptsize] at (104,3.5) {$30^\circ$};
\draw[<->, muted] (0,-6)--(126,-6); \node[below, font=\small] at (63,-6) {126 m};
\end{tikzpicture}\end{center}

\ex{2} Una escala de bombers de 10 m de longitud s'ha fixat en un punt d'un carrer. Si es recolza sobre una de les façanes, forma amb el terra un angle de $45^\circ$, i si es recolza sobre l'altra façana, l'angle és de $30^\circ$. Quina és l'amplària del carrer?

\ex{2} Una escala per a accedir a un túnel té la forma i les dimensions de la figura. Calcula la profunditat a la qual està el punt $B$ respecte del punt $A$.\begin{center}\begin{tikzpicture}[scale=0.11]
\coordinate (A) at (0,0); \coordinate (P) at ($(A)+(-30:25)$); \coordinate (Q) at ($(P)+(10,0)$); \coordinate (B) at ($(Q)+(-50:30)$);
\draw[navy, line width=1.6pt, line join=round] (A)--(P)--(Q)--(B);
\draw[muted, dashed] (P)--++(-7,0) (B)--++(-8,0);
\node[above left] at (A) {$A$}; \node[below right] at (B) {$B$};
\node[above right, font=\small] at ($(A)!0.5!(P)$) {25 m}; \node[above, font=\small] at ($(P)!0.5!(Q)$) {10 m}; \node[right, font=\small] at ($(Q)!0.5!(B)+(0.6,0)$) {30 m};
\draw ($(P)+(-4,0)$) arc (180:150:4); \node[font=\scriptsize] at ($(P)+(-6.5,1.4)$) {$30^\circ$};
\draw ($(B)+(-4,0)$) arc (180:130:4); \node[font=\scriptsize] at ($(B)+(-6.5,2.4)$) {$50^\circ$};
\end{tikzpicture}\end{center}

\ex{2} Polígons regulars:\begin{enumerate}[label=\alph*),leftmargin=*,itemsep=2pt]\item Quant mesura l'apotema d'un pentàgon regular de 10 cm de costat?\item Calcula l'àrea d'un pentàgon regular inscrit en una circumferència de 6 cm de radi.\end{enumerate}

\ex{2} Triangles isòsceles:\begin{enumerate}[label=\alph*),leftmargin=*,itemsep=2pt]\item La base d'un triangle isòsceles mesura 5 cm i l'angle oposat a aquesta base és de $55^\circ$. Calcula l'altura i l'àrea del triangle.\item El costat desigual d'un triangle isòsceles fa 34 cm i l'angle oposat, $50^\circ$. Calcula el perímetre del triangle i la seua àrea.\end{enumerate}

\ex{2} Rombes:\begin{enumerate}[label=\alph*),leftmargin=*,itemsep=2pt]\item Calcula l'àrea d'un rombe de 12 cm de costat sabent que un dels seus angles mesura $43^\circ$.\item Troba els angles d'un rombe les diagonals del qual mesuren 12 cm i 8 cm. Calcula'n també el perímetre i l'àrea.\end{enumerate}

\ex{3} Calcula la longitud de la diagonal menor i l'àrea d'un paral·lelogram de costats 10 cm i 15 cm, sabent que un dels seus angles mesura $35^\circ$.

\ex{3} Per construir un túnel d'$A$ fins a $B$ es localitza una roca $C$ visible des dels dos punts. Es mesuren les distàncies $\overline{AC}=245$ m i $\overline{BC}=658$ m, i l'angle $\widehat{ACB}=57^\circ$. Quina serà la longitud del túnel?

\ex{3} Calcula l'angle de tir $\alpha$ del jugador que està en el punt $A$, al centre de la línia de banda, si la porteria $BC$ fa 5 m i està centrada en la línia de fons.\begin{center}\begin{tikzpicture}[scale=0.055]
\fill[teal!15] (0,0) rectangle (100,60); \draw[navy, line width=1pt] (0,0) rectangle (100,60);
\draw[navy!50] (50,0)--(50,60) (50,30) circle (9);
\coordinate (A) at (50,0); \coordinate (B) at (100,27.5); \coordinate (C) at (100,32.5);
\draw[accent, line width=0.9pt] (A)--(B) (A)--(C); \fill (A) circle (0.9) (B) circle (0.9) (C) circle (0.9);
\node[below] at (A) {$A$}; \node[right] at (B) {$B$}; \node[right] at (C) {$C$};
\node[font=\scriptsize] at (84,23) {$\alpha$};
\draw[<->, muted] (0,64)--(100,64); \node[above, font=\small] at (50,64) {100 m};
\draw[<->, muted] (-4,0)--(-4,60); \node[left, font=\small] at (-4,30) {60 m};
\end{tikzpicture}\end{center}


% =====================================================
\seccio{7}{Raons trigonomètriques d'un angle qualsevol}

\exf{1} Utilitza la calculadora per a calcular les raons trigonomètriques d'aquests angles (amb quatre decimals) i indica en quin quadrant està cadascun:\begin{center}\renewcommand{\arraystretch}{1.35}\begin{tabular}{|c|C{2.2cm}|C{2.2cm}|C{2.2cm}|C{2cm}|}\hline $\alpha$ & $\sin\alpha$ & $\cos\alpha$ & $\tg\alpha$ & \textbf{Quadrant}\\\hline $120^\circ$ & & & & \\\hline $220^\circ$ & & & & \\\hline $135^\circ$ & & & & \\\hline $300^\circ$ & & & & \\\hline $240^\circ$ & & & & \\\hline $75^\circ$ & & & & \\\hline \end{tabular}\end{center}

\ex{1} Respon:\begin{tasks}(1)
  \task Si $\sin\alpha=-0{,}8$, en quin quadrant pot estar l'angle $\alpha$?
  \task Si $\cos\beta=-0{,}3$, en quin quadrant pot estar l'angle $\beta$?
  \task Si $\tg\gamma=-1{,}5$, en quin quadrant pot estar l'angle $\gamma$?
\end{tasks}

\ex{1} Calcula sense calculadora:\begin{tasks}(2)
  \task $3\sin270^\circ-2\cos180^\circ$
  \task $\cos0^\circ+\sin180^\circ-\tg180^\circ$
  \task $\sin90^\circ\cdot\cos360^\circ$
  \task $2\cos90^\circ-\sin270^\circ$
\end{tasks}

\ex{2} Calcula les altres raons trigonomètriques de l'angle en cada cas:\begin{tasks}(2)
  \task $\cos\alpha=0{,}7$ \ i \ $\alpha$ en el primer quadrant.
  \task $\tg\beta=-3$ \ i \ $\beta$ en el segon quadrant.
  \task $\sin\gamma=-\frac35$ \ i \ $\gamma$ en el tercer quadrant.
  \task $\cos\delta=\frac{5}{13}$ \ i \ $\delta$ en el quart quadrant.
\end{tasks}

\ex{2} Calcula les raons trigonomètriques de $\alpha$ sabent que $\sin\alpha=0{,}34$. Quants casos hi ha?

\ex{2} Troba els dos angles de $[0^\circ,360^\circ)$ que compleixen cada condició:\begin{tasks}(4)
  \task $\sin x=-\frac12$
  \task $\tg x=-1$
  \task $\cos x=0{,}8$
  \task $\sin x=0{,}9$
\end{tasks}\recorda{la calculadora dona un sol angle. L'altre és $180^\circ-x$ (sinus), $360^\circ-x$ (cosinus) o $x+180^\circ$ (tangent).}

\ex{2} Un angle $\beta$ tal que $0^\circ<\beta<360^\circ$ verifica $\sin\beta=-\sin30^\circ$ i $\cos\beta=\cos30^\circ$. A quin quadrant pertany? Quant mesura?

\ex{1} A quin angle entre $0^\circ$ i $360^\circ$ correspon cadascun d'aquests angles?\begin{tasks}(4)
  \task $1750^\circ$
  \task $1920^\circ$
  \task $-800^\circ$
  \task $3825^\circ$
\end{tasks}


% =====================================================
\seccio{8}{Reducció al primer quadrant}

\exf{1} Fes servir la relació entre els angles per a calcular, de manera exacta:\begin{tasks}(1)
  \task Les raons de $135^\circ$, sabent les de $45^\circ$.
  \task Les raons de $240^\circ$, sabent les de $60^\circ$.
  \task Les raons de $-30^\circ$, sabent les de $30^\circ$.
\end{tasks}

\ex{2} Calcula les raons trigonomètriques (sinus, cosinus i tangent) dels angles següents, relacionant-les amb les d'un angle del primer quadrant. Si l'angle del primer quadrant és $30^\circ$, $45^\circ$ o $60^\circ$, dona'n el valor exacte.\begin{tasks}(6)
  \task $155^\circ$
  \task $240^\circ$
  \task $390^\circ$
  \task $225^\circ$
  \task $1920^\circ$
  \task $-30^\circ$
\end{tasks}

\ex{2} Expressa les raons trigonomètriques dels angles següents en funció de les d'un angle del primer quadrant. Compte amb els signes!\begin{tasks}(5)
  \task $-40^\circ$
  \task $385^\circ$
  \task $-185^\circ$
  \task $835^\circ$
  \task $934^\circ$
\end{tasks}

\ex{2} Sabent que $\sin35^\circ\approx0{,}57$, $\cos35^\circ\approx0{,}82$ i $\tg35^\circ\approx0{,}70$, calcula sense calculadora les raons trigonomètriques de:\begin{tasks}(6)
  \task $145^\circ$
  \task $215^\circ$
  \task $325^\circ$
  \task $55^\circ$
  \task $-35^\circ$
  \task $395^\circ$
\end{tasks}

\ex{3} Expressa les raons trigonomètriques dels angles següents en funció de les d'un angle $\alpha$ tal que $0^\circ\le\alpha\le45^\circ$:\begin{tasks}(4)
  \task $120^\circ$
  \task $430^\circ$
  \task $-800^\circ$
  \task $180^\circ$
\end{tasks}\recorda{si l'angle del primer quadrant és major de $45^\circ$, passa al complementari: $\sin(90^\circ-\alpha)=\cos\alpha$, $\cos(90^\circ-\alpha)=\sin\alpha$, $\tg(90^\circ-\alpha)=\dfrac{1}{\tg\alpha}$.}

\ex{3} Simplifica: \ $\sin(180^\circ-\alpha)+\sin(180^\circ+\alpha)+\cos(360^\circ-\alpha)+\cos(180^\circ-\alpha)$.

\vfill
\begin{center}\color{muted}\small\faCheckCircle[regular]\ \ Ànim! Comprova els teus resultats amb les eines interactives dels apunts de la unitat.\end{center}

\end{document}
